\section{About the Committee}

\lead{The International Criminal Police Organization, known as INTERPOL, is the largest police cooperation body in the world, connecting the criminal police authorities of 196 member countries through a shared infrastructure of databases, secure communications, and international alerts.\source{interpolWhatIs}}

It is also one of the most misunderstood institutions in international affairs, routinely imagined as a force of cross-border detectives able to hunt and arrest suspects anywhere on earth. The reality is narrower and, for the work of this committee, far more interesting. INTERPOL commands no officers, makes no arrests, and runs no investigations of its own. It is a network that allows sovereign states to cooperate, and its authority ends precisely where their consent does. Grasping that distinction is the precondition for any serious debate on financial sovereignty, offshore secrecy, and the policing of illicit money.

\subsection{Mandate}

INTERPOL rests on the legal foundation laid down in its 1956 Constitution, where Article~2 gives the Organization two aims: to ensure the widest possible mutual assistance between all criminal police authorities, within the limits of national law and in the spirit of the Universal Declaration of Human Rights, and to develop the institutions that help prevent and suppress ordinary law crimes.\source[art.~2]{interpolConstitution} Much of what the Organization can and cannot do is already contained in that wording. The clause ``within the limits of the laws existing in the different countries'' subordinates everything it does to national sovereignty, while the reference to ``ordinary law crimes'' deliberately leaves out the political, a boundary reinforced by Article~3, which strictly forbids any intervention of a political, military, religious, or racial character.\source[art.~3]{interpolConstitution} From these provisions flow the four principles that guide INTERPOL in practice: national sovereignty, respect for human rights, neutrality, and constant and active cooperation.\source{interpolLegalDocs}

In practice INTERPOL is a network, not a police force. It works through three layers:

\begin{itemize}
  \item a General Assembly of all member states as its supreme body, an Executive Committee that oversees it between sessions, and a permanent General Secretariat in Lyon that runs daily operations;\source[arts.~5--32]{interpolConstitution}
  \item a National Central Bureau in each member state, the single link between that country's police and the network, communicating through the secure I-24/7 system;\source{interpolKeyDates}
  \item central databases covering travel documents, fingerprints, DNA, facial images, stolen vehicles, and wanted persons, plus eight colour-coded notices, of which the Red Notice, a request to locate and provisionally detain a person pending extradition, is the best known.\source{gir2026interpol}
\end{itemize}

Most relevant for us, the INTERPOL Financial Crime and Anti-Corruption Centre (IFCACC), created in 2022, coordinates the global police response to fraud, money laundering and asset recovery, and corruption, working alongside bodies such as the Financial Action Task Force and the Egmont Group of financial intelligence units.\source{interpolIFCACC2022,interpolFinancialRole}

\begin{keypoints}[What INTERPOL cannot do]
  \item It has no agents and no power of arrest, since every arrest is made by national police under national law.
  \item A Red Notice is not an international arrest warrant, but a request that each member may act on or decline, and many decline.\source{gir2026interpol}
  \item It does not investigate, prosecute, or determine guilt, all of which are matters for national and international courts.\source[p.~4]{interpolRepository2024}
  \item It is not an organ of the United Nations, but an independent intergovernmental body.
  \item Its Article~3 neutrality is contested in practice, as some states have misused notices to pursue dissidents and opponents under the cover of fraud or corruption charges, forcing INTERPOL to strengthen oversight through the Commission for the Control of INTERPOL's Files.\source{pace2017}
\end{keypoints}

This is the paradox at the heart of our agenda: the same instruments that let states recover stolen assets can be turned into tools of transnational repression, and a committee on financial sovereignty has to keep both in view at once.
